\documentclass[11pt,a4paper]{article}
\usepackage[utf8]{inputenc}
\usepackage[T1]{fontenc}
\usepackage[polish]{babel}
\usepackage{graphicx}
\usepackage{booktabs}
\usepackage{hyperref}
\usepackage{amsmath}
\usepackage{geometry}
\usepackage{pgfplots}
\pgfplotsset{compat=1.18}
\geometry{margin=2.5cm}

\title{Wykrywanie fake news -- porównanie modeli językowych\\
\large LIAR dataset: RoBERTa fine-tuning vs.\ GPT-4o zero-shot i few-shot}

\date{Czerwiec 2026}

\begin{document}
\maketitle

\section{Wprowadzenie}

Celem projektu jest porównanie różnych podejść do automatycznego wykrywania fałszywych
informacji w wypowiedziach politycznych. Zadanie realizowane jest na zbiorze LIAR
(Wang, 2017), zawierającym 12 791 krótkich wypowiedzi oznaczonych jednym z sześciu
poziomów prawdziwości przez serwis fact-checkingowy PolitiFact. Porównujemy dwa
zasadniczo różne paradygmaty: \emph{fine-tuning} pretrenowanego enkodera (RoBERTa)
oraz \emph{prompting} dużego modelu generatywnego (GPT-4o-mini i GPT-4o) w trybach
zero-shot i few-shot.

Zadanie rozważamy w dwóch wariantach:
\begin{itemize}
  \item \textbf{Klasyfikacja binarna (2kl)} -- etykiety \texttt{fake} (oryginalnie
        \texttt{pants-fire}, \texttt{false}, \texttt{barely-true}) i \texttt{real}
        (\texttt{half-true}, \texttt{mostly-true}, \texttt{true}).
  \item \textbf{Klasyfikacja sześcioklasowa (6kl)} -- pełne etykiety LIAR.
\end{itemize}

\section{Zbiór danych}

LIAR składa się z 10 240 wypowiedzi treningowych, 1 284 walidacyjnych i 1 267
testowych. Każda wypowiedź zawiera: treść, mówcę, partię polityczną, kontekst
wypowiedzi, temat oraz historię wcześniejszych ocen prawdziwości danego mówcy.

Dziesięć najczęściej omawianych tematów w zbiorze danych to: gospodarka,
opieka zdrowotna, podatki, budżet federalny, edukacja, zatrudnienie, budżet
stanowy, biografie kandydatów, wybory i imigracja\footnote{Źródło:
\url{https://aclanthology.org/P17-2067/} (Wang 2017).}.

\subsection{Rozkład klas}

W wariancie binarnym zbiór jest umiarkowanie niezbilansowany (5 752 \texttt{real}
vs 4 488 \texttt{fake} w zbiorze treningowym -- 56\% vs 44\%). W wariancie
sześcioklasowym dysproporcje są wyraźniejsze: najrzadsza klasa \texttt{pants-fire}
stanowi 8\% danych, najczęstsza \texttt{half-true} -- 21\%. Brak zbalansowania
jest istotny dla wyboru funkcji straty (sekcja \ref{sec:methods}).

\subsection{Przygotowanie tekstu}

Każda wypowiedź reprezentowana jest jako tekst łączący wiele cech:
\begin{verbatim}
{speaker} ({party}, past lies: {n}, past truths: {m}):
  {statement} Subject: {subject} Context: {context}
\end{verbatim}
Tak skonstruowany ciąg daje modelowi dostęp nie tylko do treści wypowiedzi, ale
również do meta-danych mówcy. Tekst jest tokenizowany tokenizatorem RoBERTa
z długością maksymalną 128 tokenów -- 99{,}96\% wypowiedzi mieści się w tym
limicie, mediana długości wynosi 22 tokeny.

BERT i RoBERTa jako modele pretrenowane przyjmują wejście wyłącznie w postaci
wektorów liczb całkowitych, gdzie każda wartość reprezentuje token. Konwersję
tę wykonują dedykowane tokenizatory (osobne dla BERT-a i RoBERTy), które
dodają tokeny specjalne początku i końca sekwencji, dzielą tekst na tokeny
(słowa, sub-słowa lub znaki) i przypisują każdemu unikalne ID, tworząc
wymagany wektor liczb całkowitych.

\section{Metody}\label{sec:methods}

\subsection{RoBERTa fine-tuning}

Wybraliśmy model \texttt{roberta-large} (355M parametrów) zamiast \texttt{bert-base}
z następujących powodów (Liu et al., 2019):
\begin{enumerate}
  \item \textbf{Dynamiczne maskowanie} -- RoBERTa losuje maskowane tokeny przy
        każdym przejściu, BERT stosuje statyczne maskowanie. Lepsze pokrycie
        kontekstu w trakcie pretreningu.
  \item \textbf{Większy korpus} -- 160 GB tekstu (Wikipedia, książki, CommonCrawl,
        OpenWebText, news) zamiast 16 GB BERT-a. Bogatsza wiedza ogólna,
        szczególnie istotna w domenie newsów politycznych.
  \item \textbf{Brak zadania NSP} (Next Sentence Prediction) -- udowodniono, że
        NSP nie pomaga w klasyfikacji pojedynczych zdań, a wręcz pogarsza wyniki.
  \item \textbf{Lepsze hiperparametry pretreningu} -- większe batche (8K), dłuższy
        trening.
\end{enumerate}

BERT był pretrenowany na BooksCorpus i angielskiej Wikipedii, natomiast
RoBERTa na BooksCorpus, CC-News, OpenWebText i Stories -- znacznie
większym i bardziej zróżnicowanym korpusie.

RoBERTa jest modelem typu \emph{encoder-only} -- zawiera wyłącznie część
enkoderową architektury Transformer, dzięki czemu produkuje gęsty wektor
reprezentacji wejścia (a nie generuje tekstu), co czyni ją naturalnym
wyborem do zadań klasyfikacji.

Model jest trenowany z głową klasyfikacji (Dropout $\to$ Linear $\to$ tanh
$\to$ Dropout $\to$ Linear) podpinaną do reprezentacji tokenu \texttt{[CLS]}.

\subsubsection{Hiperparametry treningu}

\begin{itemize}
  \item Liczba epok: 5
  \item Optymalizator: AdamW, $\text{lr}=1\!\cdot\!10^{-5}$, $\text{weight\_decay}=0{,}01$
  \item Batch size: 16 (A100) lub 8 (T4)
  \item Learning rate schedule: liniowy warm-up przez 1 epokę, następnie liniowy
        decay do zera (\texttt{get\_linear\_schedule\_with\_warmup})
  \item Gradient clipping: 1{,}0
  \item Selekcja epoki: najlepszy walidacyjny macro-F1 (model najlepszej epoki
        zapisywany jest na dysk)
\end{itemize}

\subsubsection{Wagi klas}

Ze względu na niezbilansowanie klas (zwłaszcza w wariancie 6kl) stosujemy
ważoną funkcję straty Cross-Entropy:
\[
w_c = \frac{N}{C \cdot n_c},
\]
gdzie $N$ to liczba wszystkich przykładów, $C$ to liczba klas, a $n_c$ to
liczba przykładów klasy $c$. Wagi te dla 2kl wynoszą $w_{\text{fake}}=1{,}14$,
$w_{\text{real}}=0{,}89$, a dla 6kl rozciągają się od ok. 0{,}81 dla
\texttt{half-true} do ok. 2{,}03 dla \texttt{pants-fire}. Mechanizm ten
przeciwdziała tendencji modelu do ignorowania klas mniejszościowych.

\subsection{GPT prompting -- zero-shot i few-shot}

Drugie podejście wykorzystuje modele GPT (\texttt{gpt-4o-mini}, \texttt{gpt-4o})
przez API OpenAI w dwóch trybach:
\begin{itemize}
  \item \textbf{Zero-shot} -- model otrzymuje wyłącznie opis zadania i jedną
        wypowiedź do oceny.
  \item \textbf{Few-shot} -- przed wypowiedzią testową prompt zawiera $K$
        przykładów z każdej klasy ze zbioru treningowego.
\end{itemize}

Przykłady few-shot dobierane są przez \emph{nearest-neighbor retrieval} oparty
na cosine similarity reprezentacji TF-IDF (1- i 2-gramy). Dla każdej klasy
wybieramy $K_{\text{per\_class}}=3$ najbardziej podobne wypowiedzi do
zapytania (łącznie 6 przykładów dla 2kl i 18 dla 6kl).

\subsection{Metryki}

Główną metryką jest \textbf{macro-F1} -- średnia F1 obliczona niezależnie dla
każdej klasy, bez ważenia liczebnością. Wybór ten jest podyktowany
niezbalansowaniem danych: accuracy oraz weighted-F1 byłyby dominowane przez
klasy większościowe i nie odzwierciedlałyby umiejętności modelu na klasach
rzadkich (krytyczne dla 6kl).

Dodatkowo raportujemy accuracy, precision i recall per klasa oraz macierze
pomyłek.

\section{Wyniki}

Rysunek \ref{fig:results_2cl} prezentuje porównanie modeli na wariancie
binarnym, a rysunek \ref{fig:results_6cl} -- na wariancie sześcioklasowym.

\begin{figure}[h]
  \centering
  \begin{tikzpicture}
    \begin{axis}[
        ybar,
        width=14cm, height=7cm,
        bar width=10pt,
        enlarge x limits=0.15,
        ymin=0, ymax=0.85,
        ylabel={Wynik na zbiorze testowym},
        symbolic x coords={RoBERTa-base, GPT-4o-mini ZS, GPT-4o-mini FS, GPT-4o ZS},
        xtick=data,
        x tick label style={font=\small, rotate=15, anchor=east},
        legend style={at={(0.5,-0.25)}, anchor=north, legend columns=3, font=\small},
        nodes near coords,
        nodes near coords style={font=\tiny},
        ymajorgrids=true, grid style=dashed,
    ]
    \addplot[fill=blue!60] coordinates {
        (RoBERTa-base, 0.640)
        (GPT-4o-mini ZS, 0.696)
        (GPT-4o-mini FS, 0.693)
        (GPT-4o ZS, 0.729)
    };
    \addplot[fill=orange!70] coordinates {
        (RoBERTa-base, 0.610)
        (GPT-4o-mini ZS, 0.695)
        (GPT-4o-mini FS, 0.686)
        (GPT-4o ZS, 0.0)
    };
    \addplot[fill=green!60!black] coordinates {
        (RoBERTa-base, 0.630)
        (GPT-4o-mini ZS, 0.697)
        (GPT-4o-mini FS, 0.692)
        (GPT-4o ZS, 0.0)
    };
    \legend{Accuracy, Macro-F1, Weighted-F1}
    \end{axis}
  \end{tikzpicture}
  \caption{Porównanie modeli na zbiorze testowym LIAR -- wariant binarny.
           Wartość 0 oznacza brak pomiaru. ZS = zero-shot, FS = few-shot (kNN).}
  \label{fig:results_2cl}
\end{figure}

\begin{figure}[h]
  \centering
  \begin{tikzpicture}
    \begin{axis}[
        ybar,
        width=12cm, height=6.5cm,
        bar width=14pt,
        enlarge x limits=0.4,
        ymin=0, ymax=0.4,
        ylabel={Wynik na zbiorze testowym},
        symbolic x coords={GPT-4o-mini ZS, GPT-4o-mini FS},
        xtick=data,
        x tick label style={font=\small},
        legend style={at={(0.5,-0.18)}, anchor=north, legend columns=3, font=\small},
        nodes near coords,
        nodes near coords style={font=\small},
        ymajorgrids=true, grid style=dashed,
    ]
    \addplot[fill=blue!60] coordinates {
        (GPT-4o-mini ZS, 0.298)
        (GPT-4o-mini FS, 0.324)
    };
    \addplot[fill=orange!70] coordinates {
        (GPT-4o-mini ZS, 0.279)
        (GPT-4o-mini FS, 0.310)
    };
    \addplot[fill=green!60!black] coordinates {
        (GPT-4o-mini ZS, 0.264)
        (GPT-4o-mini FS, 0.294)
    };
    \legend{Accuracy, Macro-F1, Weighted-F1}
    \end{axis}
  \end{tikzpicture}
  \caption{Porównanie modeli na zbiorze testowym LIAR -- wariant
           sześcioklasowy. ZS = zero-shot, FS = few-shot (kNN).}
  \label{fig:results_6cl}
\end{figure}

\subsection{Zestawienie fine-tuningu (encoder)}\label{sec:encoder_summary}

\begin{table}[h]
\centering
\caption{Porównanie modeli typu encoder fine-tuned na LIAR (zbiór testowy).
BERT-base i RoBERTa-base trenowane były tylko w wariancie 2-klasowym.}
\label{tab:encoder_summary}
\begin{tabular}{lrrrr}
\toprule
Model & 2kl Acc & 2kl Macro-F1 & 6kl Acc & 6kl Macro-F1 \\
\midrule
BERT-base-uncased                & 0{,}628 & 0{,}604 & --- & --- \\
RoBERTa-base                     & 0{,}640 & 0{,}622 & --- & --- \\
\textbf{RoBERTa-large + Optuna}  & \textbf{0{,}731} & \textbf{0{,}726} & \textbf{0{,}357} & \textbf{0{,}359} \\
\textbf{RoBERTa-large + LIAR-PLUS} & --- & --- & \textbf{0{,}383} & \textbf{0{,}380} \\
\bottomrule
\end{tabular}
\end{table}

\subsection{Najlepsze wyniki z każdej kategorii}\label{sec:best_of}

\begin{table}[h]
\centering
\caption{Najlepsze wyniki z każdej kategorii modeli (macro-F1 na zbiorze testowym).}
\label{tab:best_of}
\begin{tabular}{llrr}
\toprule
Kategoria & Model & 2kl Macro-F1 & 6kl Macro-F1 \\
\midrule
Encoder (fine-tune) & RoBERTa-large + Optuna             & 0{,}726          & \textbf{0{,}359} \\
LLM (OpenAI API)    & gpt-4o (SC=3, SO / few-shot)       & \textbf{0{,}729} & 0{,}356 \\
LLM darmowy (OpenRouter) & gpt-oss-20b (few-shot)         & 0{,}467$^{*}$    & --- \\
\bottomrule
\end{tabular}
\end{table}

W wariancie binarnym wynik jest praktycznie remisem
(RoBERTa-large 0{,}726 vs gpt-4o 0{,}729). W wariancie sześcioklasowym
RoBERTa-large nieznacznie wyprzedza GPT-4o (0{,}359 vs 0{,}356).
\\
$^{*}$ Liczone na 120-przykładowym podzbiorze testowym (limity rate-limit
darmowego OpenRouter), wynik niereprezentatywny dla pełnego zbioru.

\emph{Uwaga:} gpt-4o na 6kl został uruchomiony tylko w trybie few-shot
(zero-shot z $K=3$, $\text{SC}=3$, SO został przerwany przez błąd
\texttt{insufficient\_quota}). Na 2kl zero-shot bił few-shot (0{,}729 vs
0{,}710), więc zakładamy że analogicznie na 6kl zero-shot dałby
nieznacznie wyższe wyniki niż few-shot.

\subsection{Zestawienie wyników GPT (OpenAI API)}\label{sec:gpt_summary}

Tabele \ref{tab:gpt_2cl} i \ref{tab:gpt_6cl} podsumowują wszystkie
przeprowadzone eksperymenty GPT obok siebie. Pogrubione wiersze
wskazują najlepszy wynik w każdym wariancie.

\begin{table}[h]
\centering
\caption{Wyniki GPT na LIAR -- wariant binarny (2kl), zbiór testowy
$\sim$1267. K = liczba przykładów per klasa w few-shot, SC =
self-consistency, SO = Structured Outputs.}
\label{tab:gpt_2cl}
\begin{tabular}{llcccrr}
\toprule
Model & Tryb & K & SC & SO & Acc & Macro-F1 \\
\midrule
gpt-4o-mini & zero-shot      & -- & 1 & nie & 0{,}696 & 0{,}695 \\
gpt-4o-mini & few-shot (kNN) & 1  & 1 & nie & 0{,}693 & 0{,}686 \\
\textbf{gpt-4o} & \textbf{zero-shot} & \textbf{--} & \textbf{3} & \textbf{tak} & \textbf{0{,}729} & \textbf{0{,}729} \\
gpt-4o      & few-shot (kNN) & 3  & 3 & tak & 0{,}710 & 0{,}709 \\
\bottomrule
\end{tabular}
\end{table}

\begin{table}[h]
\centering
\caption{Wyniki GPT na LIAR -- wariant sześcioklasowy (6kl).}
\label{tab:gpt_6cl}
\begin{tabular}{llcccrr}
\toprule
Model & Tryb & K & SC & SO & Acc & Macro-F1 \\
\midrule
gpt-4o-mini & zero-shot      & -- & 1 & nie & 0{,}298 & 0{,}279 \\
gpt-4o-mini & few-shot (kNN) & 1  & 1 & nie & 0{,}324 & 0{,}310 \\
\textbf{gpt-4o} & \textbf{few-shot (kNN)} & \textbf{1}  & \textbf{1} & \textbf{tak} & \textbf{0{,}344} & \textbf{0{,}356} \\
\bottomrule
\end{tabular}
\end{table}

Najlepszy 2kl: \texttt{gpt-4o zero-shot} z self-consistency 3 i
Structured Outputs (acc=0{,}729, macro-F1=0{,}729, zero niesparowanych
odpowiedzi). Najlepszy 6kl: \texttt{gpt-4o few-shot kNN} z $K=1$
i $\text{SC}=1$ (acc=0{,}344, macro-F1=0{,}356). Eksperyment
\texttt{gpt-4o 6kl zero-shot} z próby 2 ($K=3$, $\text{SC}=3$, SO)
nie został doliczony -- wszystkie 1267 odpowiedzi zwróciły błąd
\texttt{insufficient\_quota} z OpenAI API.

\subsection{Szczegółowe wyniki GPT-4o-mini (1 próba)}

Poniższe metryki pochodzą z pojedynczego uruchomienia całego pipeline'u
(wszystkie cztery eksperymenty w jednej komórce, jedno wywołanie API).

\textbf{2kl zero-shot} (acc=0{,}696, macro-F1=0{,}695, unparsed=1, support=1266):
\begin{center}
\begin{tabular}{lrrrr}
\toprule
 & precision & recall & f1-score & support \\
\midrule
fake & 0{,}63 & 0{,}75 & 0{,}68 & 553 \\
real & 0{,}77 & 0{,}66 & 0{,}71 & 713 \\
\bottomrule
\end{tabular}
\end{center}
Macierz pomyłek: fake$\to$fake 413, fake$\to$real 140; real$\to$fake 245,
real$\to$real 468.

\textbf{2kl few-shot (kNN)} (acc=0{,}693, macro-F1=0{,}686, unparsed=18, support=1249):
\begin{center}
\begin{tabular}{lrrrr}
\toprule
 & precision & recall & f1-score & support \\
\midrule
fake & 0{,}65 & 0{,}63 & 0{,}64 & 541 \\
real & 0{,}72 & 0{,}74 & 0{,}73 & 708 \\
\bottomrule
\end{tabular}
\end{center}
Macierz pomyłek: fake$\to$fake 340, fake$\to$real 201; real$\to$fake 183,
real$\to$real 525.

\textbf{6kl zero-shot} (acc=0{,}298, macro-F1=0{,}279, unparsed=141, support=1126):
\begin{center}
\begin{tabular}{lrrrr}
\toprule
 & precision & recall & f1-score & support \\
\midrule
pants-fire   & 0{,}53 & 0{,}41 & 0{,}47 &  82 \\
false        & 0{,}25 & 0{,}62 & 0{,}36 & 218 \\
barely-true  & 0{,}15 & 0{,}01 & 0{,}02 & 187 \\
half-true    & 0{,}32 & 0{,}19 & 0{,}24 & 242 \\
mostly-true  & 0{,}31 & 0{,}38 & 0{,}34 & 212 \\
true         & 0{,}34 & 0{,}19 & 0{,}25 & 185 \\
\bottomrule
\end{tabular}
\end{center}
Model praktycznie nie używa klasy \texttt{barely-true} (recall 0{,}01) i
ma tendencję do nadmiernego predykowania klasy \texttt{false} (kolumna
\texttt{false} w macierzy pomyłek dominuje dla wszystkich klas).

\textbf{6kl few-shot (kNN)} (acc=0{,}324, macro-F1=0{,}310, unparsed=801,
support=466 -- niepełne, patrz uwaga poniżej):
\begin{center}
\begin{tabular}{lrrrr}
\toprule
 & precision & recall & f1-score & support \\
\midrule
pants-fire   & 0{,}48 & 0{,}48 & 0{,}48 & 29 \\
false        & 0{,}28 & 0{,}61 & 0{,}38 & 85 \\
barely-true  & 0{,}35 & 0{,}07 & 0{,}12 & 84 \\
half-true    & 0{,}35 & 0{,}36 & 0{,}35 & 98 \\
mostly-true  & 0{,}33 & 0{,}38 & 0{,}35 & 93 \\
true         & 0{,}35 & 0{,}12 & 0{,}17 & 77 \\
\bottomrule
\end{tabular}
\end{center}
\textit{Uwaga:} 801 z 1267 próbek nie zostało sklasyfikowanych z powodu
wyczerpania dziennego limitu API (RPD = 10\,000) -- raportowane metryki są
policzone na 466 próbkach i nie są bezpośrednio porównywalne z pozostałymi
eksperymentami.

\subsection{Próba 2 -- GPT-4o (k=3, self-consistency=3, Structured Outputs)}

Druga próba użyła silniejszego modelu \texttt{gpt-4o} z trzema rozszerzeniami:
\textbf{k=3} (3 przykłady na klasę w few-shot zamiast 1), \textbf{self-consistency=3}
(każde zapytanie powtarzane 3 razy, decyzja przez głosowanie większościowe)
oraz \textbf{Structured Outputs} (wymuszenie odpowiedzi w schemacie JSON
ze zdefiniowanym enum etykiet). Stąd \texttt{unparsed=0} dla zero-shot --
parser nigdy nie zawodzi.

\textbf{2kl zero-shot} (acc=0{,}729, macro-F1=0{,}729, unparsed=0, support=1267):
\begin{center}
\begin{tabular}{lrrrr}
\toprule
 & precision & recall & f1-score & support \\
\midrule
fake & 0{,}66 & 0{,}78 & 0{,}71 & 553 \\
real & 0{,}80 & 0{,}69 & 0{,}74 & 714 \\
\bottomrule
\end{tabular}
\end{center}
Macierz pomyłek: fake$\to$fake 430, fake$\to$real 123; real$\to$fake 220,
real$\to$real 494. Czas: 160{,}5 s (mimo $3\times$ więcej zapytań przez
self-consistency).

\textbf{2kl few-shot (kNN)} (acc=0{,}710, macro-F1=0{,}709, unparsed=795,
support=472 -- niepełne):
\begin{center}
\begin{tabular}{lrrrr}
\toprule
 & precision & recall & f1-score & support \\
\midrule
fake & 0{,}64 & 0{,}76 & 0{,}69 & 205 \\
real & 0{,}79 & 0{,}67 & 0{,}72 & 267 \\
\bottomrule
\end{tabular}
\end{center}
Macierz pomyłek: fake$\to$fake 156, fake$\to$real 49; real$\to$fake 88,
real$\to$real 179. Czas: 8489 s ($\sim$2 h 21 min) -- również wyczerpany
dzienny limit API (795 z 1267 próbek bez odpowiedzi).

\textit{Obserwacje:} GPT-4o zero-shot bije GPT-4o-mini o $\sim$3 pkt
(0{,}729 vs 0{,}696). Few-shot z 9 przykładami (3$\times$3) generuje
jeszcze dłuższe prompty niż w próbie 1, stąd jeszcze szybsze osiągnięcie
limitu dziennego.

\textbf{6kl few-shot (kNN, K=1, SC=1)} (acc=0{,}345, macro-F1=0{,}359,
unparsed=164, support=1103):
\begin{center}
\begin{tabular}{lrrrr}
\toprule
 & precision & recall & f1-score & support \\
\midrule
pants-fire   & 0{,}72 & 0{,}37 & 0{,}49 &  78 \\
false        & 0{,}35 & 0{,}48 & 0{,}40 & 218 \\
barely-true  & 0{,}27 & 0{,}40 & 0{,}32 & 185 \\
half-true    & 0{,}28 & 0{,}34 & 0{,}31 & 236 \\
mostly-true  & 0{,}46 & 0{,}15 & 0{,}23 & 213 \\
true         & 0{,}50 & 0{,}34 & 0{,}40 & 173 \\
\bottomrule
\end{tabular}
\end{center}
Czas: 2007 s ($\sim$33 min). \textit{Uwaga:} ze względu na ograniczony
budżet API uruchomiono z $K_{\text{per\_class}}=1$ oraz wyłączoną
self-consistency ($\text{SC}=1$) -- prostsza konfiguracja niż w
pozostałych eksperymentach próby 2. Mimo to model bije GPT-4o-mini 6kl
few-shot z próby 1 (0{,}359 vs 0{,}310 macro-F1) i osiąga najlepszy
wynik 6kl spośród wszystkich modeli GPT. Najtrafniej rozpoznawana jest
klasa \texttt{pants-fire} (F1=0{,}49) -- jaskrawe kłamstwa są najłatwiej
identyfikowalne.

\subsection{RoBERTa-large + LIAR-PLUS (justification)}\label{sec:liarplus}

W celu poprawy wyników w wariancie 6-klasowym zastosowaliśmy rozszerzoną wersję datasetu --
LIAR-PLUS \cite{alhindi2018} -- która dodaje do każdej wypowiedzi pole
\texttt{justification} zawierające kilka zdań uzasadnienia napisane przez
fact-checkerów PolitiFact. Tekst wejściowy modelu został wzbogacony o
to pole, długość maksymalna tokenizacji zwiększona ze 128 do 256 tokenów.
Hiperparametry: lr=5{,}79e-6, weight\_decay=0{,}124, batch=8, 7 epok
(z best params z Optuny~2kl).

\textbf{6kl + LIAR-PLUS} (acc=0{,}383, balanced acc=0{,}383,
macro-F1=0{,}380, ROC AUC macro=0{,}772):
\begin{center}
\begin{tabular}{lrrrr}
\toprule
 & precision & recall & f1-score & support \\
\midrule
pants-fire   & 0{,}40 & 0{,}40 & 0{,}40 &  92 \\
false        & 0{,}44 & 0{,}49 & 0{,}46 & 249 \\
barely-true  & 0{,}36 & 0{,}26 & 0{,}30 & 212 \\
half-true    & 0{,}35 & 0{,}30 & 0{,}32 & 265 \\
mostly-true  & 0{,}36 & 0{,}45 & 0{,}40 & 241 \\
true         & 0{,}39 & 0{,}39 & 0{,}39 & 208 \\
\bottomrule
\end{tabular}
\end{center}

Dodanie pola \texttt{justification} dało skok macro-F1 z 0{,}357 do
0{,}380 (+2{,}3 pkt) oraz accuracy z 0{,}357 do 0{,}383 (+2{,}6 pkt). Model
przebił barierę 0{,}32 z literatury, osiągając wynik porównywalny z
najnowszymi pracami na LIAR-PLUS. Co istotne, model przestał dominować
predykcji nad pojedynczymi klasami -- recall jest teraz bardziej
zrównoważony (0{,}26--0{,}49 dla każdej klasy), zamiast skupiać się na
2-3 najczęstszych.

\subsection{Open-source LLM przez OpenRouter}\label{sec:openrouter}

Jako dodatkowy baseline przetestowano darmowe modele open-source przez
platformę \href{https://openrouter.ai}{OpenRouter}: \texttt{gpt-oss-20b},
\texttt{gemma-4-31b-it}, \texttt{dolphin-mistral-24b-venice-edition},
\texttt{llama-3.3-70b-instruct}. Wszystkie to dekoder-only Transformery
różnej wielkości (20--70 mld parametrów).

\paragraph{Realne wyniki tylko dla gpt-oss-20b.} Limity dziennego użycia
darmowego planu OpenRouter (50 zapytań/dzień bez doładowania) przerwały
runy gemmy, dolphina i llamy -- wszystkie zwracały \texttt{ERROR\_STATUS\_429}
i fallback parsera przypisywał odpowiedzi do klasy \texttt{fake},
dając niereprezentatywne 0{,}442/0{,}306. Tylko gpt-oss-20b zdążył odpowiedzieć
na pełne 120 próbek przed limitem.

\begin{table}[h]
\centering
\caption{Wyniki gpt-oss-20b (OpenRouter) na 120-przykładowym podzbiorze
testowym LIAR (2kl).}
\label{tab:openrouter}
\begin{tabular}{lrrr}
\toprule
Tryb & Accuracy & Macro-F1 & Unparsed \\
\midrule
zero-shot & 0{,}467 & 0{,}407 & 48 / 120 \\
few-shot  & 0{,}500 & 0{,}467 &  0 / 120 \\
\bottomrule
\end{tabular}
\end{table}

Słabsze wyniki niż gpt-4o wynikają z trzech czynników: (1) mniejszy model
(20B vs $\sim$200B), (2) brak Structured Outputs -- parser musi szukać
\texttt{fake}/\texttt{true} w surowym tekście, (3) tryb \emph{reasoning}
modelu gpt-oss generuje wewnętrzny łańcuch rozumowania, który przy
\texttt{max\_tokens=8} bywa ucinany przed właściwą etykietą.

\subsection{Porównanie gpt-4o vs gpt-4o-mini}\label{sec:gpt4o_compare}

\begin{table}[h]
\centering
\caption{Architektura i cennik dwóch modeli OpenAI używanych w
eksperymentach. Wielkość i architektura niepublikowane przez OpenAI --
podane wartości to publiczne szacunki.}
\label{tab:gpt4o_vs_mini}
\begin{tabular}{lll}
\toprule
 & \textbf{gpt-4o-mini} & \textbf{gpt-4o} \\
\midrule
Wielkość       & $\sim$8 mld par. & $\sim$200 mld par. \\
Cena (input)   & \$0{,}15 / 1M tok. & \$2{,}50 / 1M tok. \\
Cena (output)  & \$0{,}60 / 1M tok. & \$10 / 1M tok. \\
Szybkość       & $\sim$2$\times$ szybszy & wolniejszy \\
Jakość         & słabszy w trudnych zadaniach & flagowy \\
\bottomrule
\end{tabular}
\end{table}

W naszym pipeline gpt-4o-mini użyty jako tani baseline (0{,}15\$/1M
tokenów to $\sim$17$\times$ taniej niż gpt-4o), gpt-4o jako model główny
do uzyskania najlepszych wyników. W praktyce gpt-4o przewyższa
gpt-4o-mini o 3--4 punkty macro-F1 na 2kl (0{,}729 vs 0{,}695) i
4--5 punkty na 6kl (0{,}356 vs 0{,}310).

\subsection{Czas wykonania}

Pomiary czasu inferencji na pełnym zbiorze testowym (1267 wypowiedzi) dla
modeli GPT przez API:

\begin{center}
\begin{tabular}{lrr}
\toprule
Eksperyment & Czas całkowity & Czas na próbkę \\
\midrule
GPT-4o-mini 2kl zero-shot   &   94 s &   75 ms \\
GPT-4o-mini 2kl few-shot    &  163 s &  130 ms \\
GPT-4o-mini 6kl zero-shot   &  714 s &  560 ms \\
GPT-4o-mini 6kl few-shot    & 4001 s & 3160 ms \\
\midrule
\textbf{Suma}               & \textbf{$\sim$83 min} & \\
\bottomrule
\end{tabular}
\end{center}

Few-shot 6-klasowy wyczerpał dzienny limit zapytań API (RPD = 10\,000),
przez co 801 z 1267 próbek zwróciło błąd zamiast predykcji
(\texttt{unparsed=801}). Raportowane metryki dla tego eksperymentu są
policzone na 466 próbkach, a nie na pełnym zbiorze testowym, i nie są
bezpośrednio porównywalne z pozostałymi. Few-shot 18-przykładowy (3 na
klasę $\times$ 6 klas) generuje długie prompty, które szybko przekraczają
limit tokenów na minutę (TPM = 200\,000).

\section{Dyskusja}

\subsection{Few-shot vs.\ zero-shot}

W obu wariantach (2kl i 6kl) prompting few-shot daje wyniki porównywalne
z zero-shot, a w binarnym przypadku nawet nieznacznie gorsze. Obserwacja ta
jest zgodna z literaturą i wynika ze specyfiki zadania: LIAR wymaga
\emph{wiedzy faktograficznej} (czy konkretne stwierdzenie polityczne jest
prawdziwe), a nie rozpoznawania wzorca stylistycznego. Przykłady w prompcie
mogą zilustrować format odpowiedzi -- co GPT i tak rozumie w trybie zero-shot
-- ale nie dostarczają wiedzy potrzebnej do oceny nowej, niezwiązanej tematycznie
wypowiedzi. Dodatkowo wyszukiwanie przykładów przez kNN po podobieństwie
tematycznym może wręcz mylić model: wkleja on do promptu wypowiedzi o tym
samym temacie, ale o różnych etykietach prawdziwości, sugerując że temat nie
niesie informacji o etykiecie.

Inaczej zachowałby się klasyfikator spamu: tam wystarczy pokazać kilka
przykładów stylu, by model nauczył się wzorca leksykalnego. W LIAR wzorca
takiego nie ma -- prawdziwość zależy od stanu świata.

\subsection{Trudność wariantu 6-klasowego}

Spadek wyniku z $\sim 0{,}70$ accuracy (2kl) do $\sim 0{,}30$ (6kl) jest
spodziewany. Granice między klasami \texttt{barely-true}, \texttt{half-true}
i \texttt{mostly-true} są w dużej mierze subiektywne -- nawet anotatorzy
PolitiFact niekoniecznie byliby zgodni co do różnicy między \texttt{half-true}
a \texttt{mostly-true}. Modele osiągają wyniki w okolicy aktualnego sufitu
literatury dla tego zadania (27--35\% accuracy).

\begin{table}[h]
\centering
\caption{Wyniki ewaluacji na zbiorze LIAR w oryginalnej publikacji Wang
2017 \cite{wang2017}. Górna sekcja: modele tekstowe. Dolna sekcja: modele
hybrydowe (tekst + metadane).}
\label{tab:wang2017}
\begin{tabular}{lrr}
\toprule
Models & Valid. & Test \\
\midrule
Majority             & 0{,}204 & 0{,}208 \\
SVMs                 & 0{,}258 & 0{,}255 \\
Logistic Regression  & 0{,}257 & 0{,}247 \\
Bi-LSTMs             & 0{,}223 & 0{,}233 \\
CNNs                 & 0{,}260 & 0{,}270 \\
\midrule
\multicolumn{3}{l}{Hybrid CNNs} \\
Text + Subject       & 0{,}263 & 0{,}235 \\
Text + Speaker       & \textbf{0{,}277} & 0{,}248 \\
Text + Job           & 0{,}270 & 0{,}258 \\
Text + State         & 0{,}246 & 0{,}256 \\
Text + Party         & 0{,}259 & 0{,}248 \\
Text + Context       & 0{,}251 & 0{,}243 \\
Text + History       & 0{,}246 & 0{,}241 \\
Text + All           & 0{,}247 & \textbf{0{,}274} \\
\bottomrule
\end{tabular}
\end{table}

Tabela~\ref{tab:wang2017} pokazuje, że nawet autorzy zbioru LIAR uzyskali
najlepszą accuracy 0{,}274 na zbiorze testowym (Hybrid CNN z pełnymi
metadanymi), a pojedyncza klasyfikacja większościowa daje 0{,}208. Słaby
wynik klasyfikacji 6-klasowej nie wynika więc z błędu modelu, lecz z
nieusuwalnej dwuznaczności semantycznej -- klasy \texttt{half-true} i
\texttt{mostly-true} wymagają zewnętrznej wiedzy o świecie i często nie są
rozróżnialne nawet dla ludzkiego anotatora. Nasze wyniki (RoBERTa-large
6kl: 0{,}357, RoBERTa-large + LIAR-PLUS 6kl: 0{,}383, GPT-4o few-shot:
0{,}344) mieszczą się w zakresie raportowanym w nowszej literaturze
na tym zadaniu.

Systematyczny przegląd literatury~\cite{emergentmind} potwierdza istnienie
\emph{performance ceiling} na poziomie $\sim$0{,}32 dla klasyfikacji
6-klasowej i $\sim$0{,}62 dla klasyfikacji binarnej:

\begin{table}[h]
\centering
\caption{Wybrane wyniki z literatury na zbiorze LIAR (źródło:
emergentmind.com/topics/liar-benchmark).}
\label{tab:liar_ceiling}
\begin{tabular}{lrr}
\toprule
Model & Test Acc & Weighted F1 \\
\midrule
Linear SVM (BoW) -- 2kl                  & 0{,}624 & 0{,}316 \\
RoBERTa (prior) -- 2kl                   & 0{,}620 & -- \\
XGBoost (GloVe, overfit) -- 6kl          & 0{,}249 & -- \\
GPT-3 Curie (fine-tuned) -- 6kl          & 0{,}295 & -- \\
CNN (text-only, Wang 17) -- 6kl          & 0{,}270 & -- \\
\midrule
\textbf{RoBERTa-large + Optuna (nasze)} -- 2kl & \textbf{0{,}738} & 0{,}735 \\
\textbf{RoBERTa-large + Optuna (nasze)} -- 6kl & \textbf{0{,}357} & 0{,}357 \\
\textbf{RoBERTa-large + LIAR-PLUS (nasze)} -- 6kl & \textbf{0{,}383} & 0{,}379 \\
\textbf{GPT-4o few-shot (nasze)} -- 6kl  & \textbf{0{,}344} & 0{,}341 \\
\bottomrule
\end{tabular}
\end{table}

Cytat za przeglądem: \emph{,,fine-grained classification accuracy and
weighted F1 do not exceed 0{,}32, even when employing linear SVM, XGBoost,
and RoBERTa-based architectures, [...] reflecting exploitation of lexical
artifacts rather than transferability''}. Co więcej, modele drzewiaste
przy treningu osiągają $\sim$99\% accuracy na zbiorze treningowym, ale
zaledwie $\sim$25\% na zbiorze testowym -- klasyczny obraz uczenia się
płytkich, leksykalnych wzorców zamiast prawdziwej wiedzy faktograficznej.
Nasze wyniki -- RoBERTa-large 2kl (0{,}738) oraz GPT-4o 6kl (0{,}344) --
sytuują się w górnej części zakresu raportowanego w literaturze dla tego
zbioru.

Co więcej, sami fact-checkerzy PolitiFact -- tworzący etykiety dla zbioru
LIAR -- regularnie nie zgadzają się ze sobą przy przypisywaniu jednej z
sześciu klas. Granice między \texttt{barely-true}, \texttt{half-true} i
\texttt{mostly-true} są określone w wewnętrznych wytycznych redakcji, ale
ich stosowanie do konkretnych wypowiedzi zależy od interpretacji
poszczególnego dziennikarza. Liczne wypowiedzi polityków bywały
re-klasyfikowane lub publicznie krytykowane jako niespójne, co dowodzi że
nawet ludzki ekspert nie osiąga pełnej zgodności na tym zadaniu. Modele
uczone na tak zaszumianych etykietach nie mogą przekroczyć tzw.
\emph{anotacyjnego sufitu} (\emph{annotation ceiling}) -- czyli górnej
granicy jakości wynikającej z subiektywności samych danych
treningowych.

\subsection{Techniki podnoszące jakość modeli}

Aby maksymalnie wycisnąć jakość mimo trudności zadania, zastosowaliśmy
następujące techniki:

\begin{itemize}
  \item \textbf{Ważona funkcja straty (class weights)} -- Cross-Entropy
        z wagami $w_c = N / (C \cdot n_c)$ przeciwdziała tendencji modelu
        do ignorowania klas mniejszościowych (zwłaszcza \texttt{pants-fire}
        w wariancie 6kl).
  \item \textbf{Label smoothing} ($\varepsilon=0{,}1$) -- redukuje
        nadmierną pewność predykcji i poprawia kalibrację, ograniczając
        eksplozję val\_loss przy fine-tuningu RoBERTy-large.
  \item \textbf{Wyszukiwanie hiperparametrów Optuna} -- 20 triali z
        samplerem TPE i pruningiem MedianPruner po pierwszej epoce.
        Optymalizowane: \texttt{lr}, \texttt{weight\_decay},
        \texttt{epochs}, \texttt{batch\_size}. Best params: lr=5{,}79e-6,
        wd=0{,}124, epochs=5, batch=8.
  \item \textbf{Early stopping} -- model najlepszej epoki (po val
        macro-F1) zapisywany na dysk; trening przerywany po 2 epokach
        bez poprawy.
  \item \textbf{Linear warmup + decay} -- jeden epoch warmup, następnie
        liniowy spadek lr do zera. Stabilizuje fine-tuning pretrenowanego
        enkodera.
  \item \textbf{Gradient clipping} ($\|\nabla\|_2 \le 1{,}0$) --
        zabezpiecza przed skokami gradientu na trudnych batchach.
  \item \textbf{Mixed precision (bfloat16)} -- przyspieszenie treningu
        $\sim$2$\times$ na A100 bez utraty jakości, dzięki czemu Optuna
        wykonała 20 triali w $\sim$2 godziny.
  \item \textbf{Structured Outputs} (GPT) -- wymuszenie odpowiedzi w
        formacie JSON ze schematem \texttt{enum} etykiet, co eliminuje
        problem niezparsowanych odpowiedzi (z $\sim$11\% do $<$1\%).
  \item \textbf{Self-consistency} (GPT, próba 2) -- $N=3$ niezależne
        zapytania per próbkę i głosowanie większościowe, redukujące
        wariancję stochastyczną modelu generatywnego.
  \item \textbf{Few-shot kNN retrieval} -- przykłady demonstracyjne
        dobierane przez TF-IDF cosine similarity, a nie losowo --
        kontekst tematycznie zbliżony do zapytania.
  \item \textbf{Selekcja po macro-F1, nie accuracy} -- zarówno selekcja
        najlepszej epoki, jak i raportowanie głównego wyniku odbywa się
        na podstawie macro-F1, zgodnie z zaleceniami literatury LIAR.
        Macro-F1 traktuje każdą klasę jednakowo, niezależnie od jej
        liczebności -- accuracy i weighted-F1 byłyby zdominowane przez
        klasy większościowe i nie odzwierciedlałyby zdolności modelu do
        rozpoznawania klas rzadkich (krytyczne dla 6kl, gdzie
        \texttt{pants-fire} stanowi tylko 8\% danych).
\end{itemize}

\subsection{Fine-tuning vs.\ prompting}

Wyniki sugerują, że fine-tuning RoBERTy-base nie wystarcza do pokonania
prostego promptingu GPT-4o-mini (0{,}64 vs 0{,}70 accuracy w wariancie 2kl).
Spodziewamy się, że \texttt{roberta-large} z odpowiednio dobranymi
hiperparametrami zbliży się do wyników GPT-4o, lecz nie powinien ich
zdominować -- granicą jest jakość samych etykiet oraz brak dostępu modelu
do zewnętrznej wiedzy o świecie. Przewaga GPT-4o nad mniejszymi modelami
jest spójna z ogólnym obserwowaniem, że większe modele wykazują lepszą
\emph{world knowledge} przydatną w zadaniach fact-checkingowych.

\section{Wnioski}

\begin{enumerate}
  \item RoBERTa-large jest naturalnym wyborem nad BERT-em ze względu na
        większy korpus pretreningu, dynamiczne maskowanie i lepsze
        właściwości reprezentacji do klasyfikacji pojedynczych zdań.
  \item W zadaniach wymagających wiedzy o świecie (jak fact-checking)
        few-shot nie poprawia istotnie wyników nad zero-shot -- mechanizm
        ten jest skuteczny tam, gdzie kluczowy jest wzorzec, nie fakty.
  \item Macro-F1 jest właściwą metryką dla niezbilansowanego LIAR-a;
        accuracy maskuje słabe wyniki na klasach mniejszościowych.
  \item Wagi klas i Structured Outputs to dwie najistotniejsze poprawki
        techniczne, które dają mierzalny zysk małym kosztem implementacyjnym.
\end{enumerate}

\section{Możliwości dalszego rozwoju}

Naturalnym kierunkiem rozwoju jest \textbf{klasyfikacja kaskadowa}:
najpierw model binarny rozstrzyga, czy wypowiedź jest \texttt{fake} czy
\texttt{real}, a następnie -- w obrębie każdej z dwóch grup -- drugi
model przypisuje jedną z trzech oryginalnych etykiet LIAR
(\texttt{pants-fire}/\texttt{false}/\texttt{barely-true} po stronie fake;
\texttt{half-true}/\texttt{mostly-true}/\texttt{true} po stronie real).
Pozwala to każdemu modelowi specjalizować się w prostszym zadaniu i
może ograniczyć błędy między bliskimi semantycznie klasami
(np. \texttt{half-true} vs \texttt{mostly-true}).

\begin{thebibliography}{9}
  \bibitem{wang2017}
  W.\ Y.\ Wang, ``Liar, liar pants on fire'': A new benchmark dataset for fake
  news detection, \emph{ACL}, 2017.

  \bibitem{liu2019}
  Y.\ Liu et al., RoBERTa: A robustly optimized BERT pretraining approach,
  \emph{arXiv:1907.11692}, 2019.

  \bibitem{devlin2019}
  J.\ Devlin et al., BERT: Pre-training of deep bidirectional transformers
  for language understanding, \emph{NAACL}, 2019.

  \bibitem{openai2024}
  OpenAI, GPT-4o system card, 2024.

  \bibitem{labellerr2024}
  Labellerr, RoBERTa: A robustly optimized BERT pretraining approach,
  \url{https://www.labellerr.com/blog/roberta-a-robustly-optimized-bert-pretraining-approach/}

  \bibitem{emergentmind}
  Emergent Mind, LIAR Benchmark for Fake News Detection,
  \url{https://www.emergentmind.com/topics/liar-benchmark}

  \bibitem{alhindi2018}
  T.\ Alhindi, S.\ Petridis, S.\ Muresan, Where is your evidence:
  Improving fact-checking by justification modeling, \emph{Proceedings of
  the First Workshop on Fact Extraction and Verification (FEVER)}, 2018.
  \url{https://github.com/Tariq60/LIAR-PLUS}
\end{thebibliography}

\end{document}
